% 第6章贯穿例子：两封邮件 x1,x2 上的六种可实现预测（标签网格视图）。
% 行 = x1 的预测，列 = x2 的预测，次序均为 正常、广告、诈骗。
% 实心点 = 限制类 V 中的模式：
%   h1=(正常,正常) h2=(正常,广告) h3=(广告,广告)
%   h4=(广告,诈骗) h5=(诈骗,诈骗) h6=(诈骗,正常)
% 每行、每列恰有两个实心点；任意两行只共享一列（无完整四角矩形）。
\begingroup
\definecolor{LtSixGreen}{HTML}{829C8B}
\begin{tikzpicture}[figure picture,x=1mm,y=1mm]
 \def\cs{12}
 % grid: column c (1..3) x = 22+\cs*(c-1); row r (1..3) y = 30-\cs*(r-1)
 \foreach \r in {1,2,3}{\foreach \c in {1,2,3}{
   \draw[FigureGrid,line width=.5pt] ({16+\cs*(\c-1)},{24-\cs*(\r-1)}) rectangle ++(\cs,\cs);}}
 \foreach \lab [count=\c] in {正常,广告,诈骗}{
   \node[anchor=south] at ({22+\cs*(\c-1)},37) {\lab};}
 \foreach \lab [count=\r] in {正常,广告,诈骗}{
   \node[anchor=east] at (15,{30-\cs*(\r-1)}) {\lab};}
 \node[figure note,anchor=south] at (34,42) {$x_2$ 上的预测};
 \node[figure note,rotate=90,anchor=south] at (4,18) {$x_1$ 上的预测};
 % empty cells
 \foreach \r/\c in {1/3,2/1,3/2}{
   \draw[FigureMuted,line width=.6pt,fill=white] ({22+\cs*(\c-1)},{30-\cs*(\r-1)}) circle (1.5);}
 % filled cells with hypothesis names
 \foreach \r/\c/\h in {1/1/1,1/2/2,2/2/3,2/3/4,3/3/5,3/1/6}{
   \fill[FigureBlue] ({22+\cs*(\c-1)},{30-\cs*(\r-1)}) circle (1.9);
   \node[font=\figurefont\scriptsize,text=FigureInk,anchor=north west,inner sep=.5pt]
     at ({22.8+\cs*(\c-1)},{28.6-\cs*(\r-1)}) {$h_{\h}$};}
 % list on the right
 \node[figure heading,anchor=west] at (60,40) {限制类 $\mathcal V=\mathcal H|_{\{x_1,x_2\}}$};
 \foreach \txt [count=\i] in {{$h_1$：正常，正常},{$h_2$：正常，广告},{$h_3$：广告，广告},
     {$h_4$：广告，诈骗},{$h_5$：诈骗，诈骗},{$h_6$：诈骗，正常}}{
   \node[anchor=west] at (62,{37.5-5.2*\i}) {\txt};}
 \fill[FigureBlue] (62,-1) circle (1.5); \node[figure note,anchor=west] at (64,-1) {可实现};
 \draw[FigureMuted,line width=.6pt,fill=white] (80,-1) circle (1.2); \node[figure note,anchor=west] at (82,-1) {不可实现};
\end{tikzpicture}
\endgroup
